% Preámbulo del Documento de Ingeniería del Proyecto – Fase 1 (Equipo B).
% Formato del machote 0.3: carta, márgenes de 2.5 cm, 11 pt, páginas numeradas.
\PassOptionsToPackage{unicode}{hyperref}
\PassOptionsToPackage{hyphens}{url}
\documentclass[11pt,letterpaper,spanish]{article}
\usepackage[margin=2.5cm]{geometry}
\usepackage[T1]{fontenc}
\usepackage[utf8]{inputenc}
\usepackage{textcomp}
\usepackage{lmodern}
\usepackage{xcolor}
\usepackage{amsmath,amssymb}
\usepackage{microtype}
\usepackage{parskip}
\usepackage{longtable,booktabs,array}
\usepackage{calc}
\newcounter{none} % tablas sin número (formato heredado de la conversión)
\usepackage{etoolbox}
\makeatletter
\patchcmd\longtable{\par}{\if@noskipsec\mbox{}\fi\par}{}{}
\makeatother
\usepackage{footnotehyper}
\makesavenoteenv{longtable}
\usepackage{graphicx}
\usepackage{float}
\usepackage[labelformat=empty,font=small]{caption}
\usepackage{soul}
\usepackage[bidi=default,shorthands=off]{babel}
\usepackage{newunicodechar}
\newunicodechar{☐}{$\square$}
\newunicodechar{✓}{\checkmark}
\newunicodechar{≤}{$\leq$}
\newunicodechar{≥}{$\geq$}
\newunicodechar{→}{$\rightarrow$}
\usepackage{fancyhdr}
\pagestyle{fancy}
\fancyhf{}
\fancyhead[L]{\small IS1 · Equipo B · Sistema de Órdenes de Servicio y Citas}
\fancyhead[R]{\small v\versiondoc}
\fancyfoot[C]{\thepage}
\setlength{\emergencystretch}{3em}
\providecommand{\tightlist}{\setlength{\itemsep}{0pt}\setlength{\parskip}{0pt}}
\setcounter{secnumdepth}{-\maxdimen} % la numeración va en el título, como en el machote
\setcounter{tocdepth}{2}
\usepackage{bookmark}
\usepackage{xurl}
\urlstyle{same}
\hypersetup{pdflang={es}, hidelinks, pdftitle={IS1-Fase1-EquipoB-v\versiondoc}}

% ---------------------------------------------------------------------------
% Figuras. El machote pide número, título y ruta del archivo fuente.
%   \figura{archivo}{opciones de tamaño}{número}{título}{ruta fuente}
\newcommand{\figura}[5]{%
  \begin{figure}[H]\centering
  \includegraphics[#2,keepaspectratio]{#1}
  \caption{\textbf{Figura #3.} #4. Fuente: \texttt{#5}}
  \end{figure}}

% Figura pendiente: se deja el lugar, la descripción y la fuente, en rojo,
% sin la imagen. Se reemplaza por \figura cuando exista la imagen.
%   \figurapendiente{número}{título}{descripción de lo que debe mostrar}{ruta fuente}
\newcommand{\figurapendiente}[4]{%
  \begin{figure}[H]\centering
  \fcolorbox{red}{red!4}{\parbox{0.92\linewidth}{\color{red}\small
    \textbf{PENDIENTE: imagen de la Figura #1}\par\smallskip
    \textbf{Debe mostrar:} #3\par\smallskip
    \textbf{Archivo fuente:} \texttt{#4}}}
  \caption{\textcolor{red}{\textbf{Figura #1.} #2. Fuente: \texttt{#4}}}
  \end{figure}}

% Texto pendiente (dato que falta o que debe confirmarse con el cliente).
\newcommand{\pendiente}[1]{\textcolor{red}{\textbf{[Pendiente:} #1\textbf{]}}}
